On écrase un comprimé de «Vitamine C 500», et on dissout la poudre dans l'eau pour obtenir une solution aqueuse ( $S_{A}$ ) de volume $V_{0}=200 m L$ et de concentration molaire $C_{A}$.\\
On dose le volume $V_{A}=20 m L$ de la solution ( $S_{A}$ ), avec une solution aqueuse d'hydroxyde de sodium $N a_{(a q)}^{+}+H O_{(a q)}^{-}$de concentration molaire $C_{B}=2,0.10^{-2} m o l . L^{-1}$. L'équivalence est obtenue pour un volume de solution versée $V_{B, E}=14,2 m L$.\\
2.1. Écrire l'équation chimique modélisant la transformation qui a eu lieu lors du dosage.\\
2.2. Calculer la concentration molaire $C_{A}$.\\
2.3. En déduire la valeur de la masse d'acide ascorbique contenu dans ce comprimé puis expliquer l'indication «Vitamine C 500». On donne : $M\left(C_{6} H_{8} O_{6}\right)=176 g$. mol $^{-1}$.


